\chapter{Speculative Optimization and Dynamic Deoptimization}
\label{chap:speculative}

\section{The Role of Speculation in SSA Supercompilation}
\label{sec:speculative:intro}

In programs utilizing dynamic types, polymorphism, or unpredictable control flow, static driving must frequently widen configurations to preserve soundness across worst-case scenarios. However, in practice, program execution overwhelmingly follows predictable monomorphic types and single-path branches.

NumLang resolves this tension via \emph{Speculative Optimization and Dynamic Deoptimization} (Phase 52, \texttt{src/mir/speculate.rs}).

\section{Type-Profile Analysis (Phase 52)}
\label{sec:speculative:profiling}

During Tier 0 JIT execution (Chapter~\ref{chap:architecture}), NumLang injects lightweight profiling hooks at function entries and call sites. If a dynamic interface call or variant pattern match receives the same concrete type tag in $>99.5\%$ of observed executions, the compiler marks the site as a \emph{speculative candidate}.

\section{The TypeGuard Terminator}
\label{sec:speculative:typeguard}

During Tier 1 supercompilation, the driver does not bifurcate on speculative sites. Instead, it emits a \texttt{Terminator::TypeGuard}:

\begin{lstlisting}[language=Rust, caption={The TypeGuard Terminator in NumLang.}]
pub enum Terminator {
    TypeGuard {
        target_place: Place,
        expected_tag: u64,
        fast_path: BasicBlockId,
        deopt_id: DeoptId,
    },
    // ...
}
\end{lstlisting}

The fast path assumes the expected tag unconditionally, allowing the supercompiler to inline, deforest, and optimize the body as if it were statically monomorphic.

\section{Deoptimization Stubs and Frame Reconstruction}
\label{sec:speculative:deopt_stubs}

If runtime execution violates the speculative assumption ($\text{tag} \ne \text{expected\_tag}$):
\begin{enumerate}
    \item The fast-path branch traps into a specialized \emph{deoptimization stub} (\texttt{src/codegen/cranelift/deopt.rs}).
    \item The stub reads the current native machine registers, decodes the \texttt{DeoptId} metadata table, and reconstructs the unoptimized interpreter stack frame.
    \item Execution resumes seamlessly inside the unoptimized Tier 0 interpreter without observable program interruption.
\end{enumerate}
This architecture enables aggressive speculative metacomputation with zero sacrifice of language safety.

\section{Speculative Type Guard Synthesis}
\label{sec:speculate:impl}

Listing~\ref{lst:speculate_impl} presents the type-profiling and speculative guard injection logic from \texttt{src/mir/speculate.rs}:

\begin{lstlisting}[language=Rust, caption={Speculative Type Guard Generation (\texttt{src/mir/speculate.rs}).}, label={lst:speculate_impl}]
//!
//! Provides runtime profiling of observed dynamic type tags at polymorphic sites,
//! computes statistical confidence scores, and inserts speculative `Terminator::TypeGuard`
//! fast-path branches in residualized MIR when confidence exceeds configurable thresholds.

use std::collections::HashMap;
use crate::mir::{BasicBlockId, Terminator};
use crate::mir::lower::{MirBasicBlock, MirFunction, MirProgram};

/// Concrete observed type profile for a local variable or parameter at a specialization site.
#[derive(Debug, Clone, PartialEq, serde::Serialize, serde::Deserialize)]
pub struct TypeProfile {
    /// Identifier of the local or parameter.
    pub local: String,
    /// Most frequently observed concrete type tag (e.g. 1 for i64, 2 for f64, 3 for ptr).
    pub observed_tag: i64,
    /// Confidence fraction in the range [0.0, 1.0] (observed_count / total_count).
    pub confidence: f64,
    /// Total number of recorded execution samples.
    pub sample_count: usize,
}

/// Dynamic type profiler tracking observations across execution runs.
#[derive(Debug, Default, Clone)]
pub struct TypeProfiler {
    observations: HashMap<String, HashMap<i64, usize>>,
}

impl TypeProfiler {
    pub fn new() -> Self {
        Self {
            observations: HashMap::new(),
        }
    }

    /// Record a single concrete type tag observation for a given local identifier.
    pub fn record_observation(&mut self, local: &str, tag: i64) {
        *self
            .observations
            .entry(local.to_string())
            .or_default()
            .entry(tag)
            .or_insert(0) += 1;
    }

    /// Compute concrete type profiles for all observed locals.
    pub fn compute_profiles(&self) -> HashMap<String, TypeProfile> {
        let mut profiles = HashMap::new();
        for (local, tag_counts) in &self.observations {
            let total: usize = tag_counts.values().sum();
            if total == 0 {
                continue;
            }
            if let Some((&dominant_tag, &count)) = tag_counts.iter().max_by_key(|(_, &c)| c) {
                let confidence = count as f64 / total as f64;
                profiles.insert(
                    local.clone(),
                    TypeProfile {
                        local: local.clone(),
                        observed_tag: dominant_tag,
                        confidence,
                        sample_count: total,
                    },
                );
            }
        }
        profiles
    }

    /// Clear all recorded observations.
    pub fn clear(&mut self) {
        self.observations.clear();
    }
}

/// Analyze a function and insert speculative type guards at polymorphic/indirect call sites
/// where the observed profile confidence meets or exceeds `confidence_threshold` (default: 0.95).
///
/// Returns the number of speculative type guards successfully inserted into the function.
pub fn insert_speculative_type_guards(
    func: &mut MirFunction,
    profiles: &HashMap<String, TypeProfile>,
    confidence_threshold: f64,
) -> usize {
    if func.blocks.is_empty() {
        return 0;
    }

    let mut guards_inserted = 0;
    let mut new_blocks: Vec<MirBasicBlock> = Vec::new();

    // Iterate through blocks and look for candidates to guard
    for block in &func.blocks {
        let mut split_occurred = false;

        if let Terminator::IndirectCall { callee, args, dest, next } = &block.terminator {
            // Check if any argument or the callee place has a high-confidence type profile
            let candidate_local = if let Some(prof) = profiles.get(&callee.local) {
                if prof.confidence >= confidence_threshold {
                    Some((callee.clone(), prof.observed_tag))
                } else {
                    None
                }
            } else {
                args.iter().find_map(|arg| {
                    profiles.get(&arg.local).and_then(|prof| {
                        if prof.confidence >= confidence_threshold {
                            Some((arg.clone(), prof.observed_tag))
                        } else {
                            None
                        }
                    })
                })
            };

            if let Some((guard_place, expected_tag)) = candidate_local {
                split_occurred = true;
                guards_inserted += 1;

                let fast_id = BasicBlockId(func.blocks.len() + new_blocks.len() + 1);
                let deopt_id = BasicBlockId(func.blocks.len() + new_blocks.len() + 2);

                // Original block now ends with TypeGuard
                let guard_block = MirBasicBlock {
                    id: block.id.clone(),
                    statements: block.statements.clone(),
                    terminator: Terminator::TypeGuard {
                        local: guard_place,
                        expected_tag,
                        fast_path: fast_id.clone(),
                        deopt_stub: deopt_id.clone(),
                    },
                    arguments: vec![],
                };
                new_blocks.push(guard_block);

                // Fast-path block: optimized monomorphic execution continuing to `next`
                let fast_block = MirBasicBlock {
                    id: fast_id,
                    statements: vec![],
                    terminator: Terminator::IndirectCall {
                        callee: callee.clone(),
                        args: args.clone(),
                        dest: dest.clone(),
                        next: next.clone(),
                    },
                    arguments: vec![],
                };
                new_blocks.push(fast_block);

                // Deopt stub block: unspecialized fallback path
                let deopt_block = MirBasicBlock {
                    id: deopt_id,
                    statements: vec![],
                    terminator: Terminator::IndirectCall {
                        callee: callee.clone(),
                        args: args.clone(),
                        dest: dest.clone(),
                        next: next.clone(),
                    },
                    arguments: vec![],
                };
                new_blocks.push(deopt_block);
            }
        }

        if !split_occurred {
            new_blocks.push(block.clone());
        }
    }

    if guards_inserted > 0 {
        func.blocks = new_blocks;
    }

    guards_inserted
}

/// Optimize an entire MirProgram with speculative type guards using profiler statistics.
pub fn speculate_program(
    program: &mut MirProgram,
    profiler: &TypeProfiler,
    confidence_threshold: f64,
) -> usize {
    let profiles = profiler.compute_profiles();
    let mut total_guards = 0;
    for func in &mut program.functions {
        total_guards += insert_speculative_type_guards(func, &profiles, confidence_threshold);
    }
    total_guards
}

/// Extract type profile statistics for a function's parameters and locals.
pub fn collect_function_type_profiles(
    func: &MirFunction,
    profiler: &TypeProfiler,
) -> HashMap<String, TypeProfile> {
    let all_profiles = profiler.compute_profiles();
    let mut func_profiles = HashMap::new();
    for (p_name, _) in &func.params {

\end{lstlisting}

\section{Interpreter Frame Reconstruction and Deoptimization}
\label{sec:speculate:deopt_impl}

Listing~\ref{lst:deopt_impl} shows the deoptimization stub lowering and frame reconstruction from \texttt{src/codegen/cranelift/deopt.rs}:

\begin{lstlisting}[language=Rust, caption={Deoptimization Stubs and Frame Reconstruction (\texttt{src/codegen/cranelift/deopt.rs}).}, label={lst:deopt_impl}]
//!
//! Provides metadata tables for reconstructing unspecialized interpreter call frames
//! when speculative type guards fail, runtime deoptimization stubs, and atomic OSR
//! transition slots for hot-path upgrade from baseline to specialized native code.

use std::collections::HashMap;
use std::sync::atomic::{AtomicPtr, AtomicU64, Ordering};
use std::sync::RwLock;
use crate::typecheck::Type;

/// Metadata describing a deoptimization safepoint site.
#[derive(Debug, Clone, serde::Serialize, serde::Deserialize)]
pub struct DeoptMetadata {
    /// Unique safepoint identifier.
    pub deopt_id: u32,
    /// Function enclosing this safepoint.
    pub func_name: String,
    /// Unspecialized basic block ID to resume execution at.
    pub resume_bb: usize,
    /// Size of the interpreter frame in bytes.
    pub frame_size: usize,
    /// Live variables at the safepoint: (variable_name, type, offset_or_slot).
    pub live_vars: Vec<(String, Type, i64)>,
}

/// Registry storing metadata for all compiled deoptimization points.
#[derive(Debug, Default)]
pub struct DeoptTable {
    entries: HashMap<u32, DeoptMetadata>,
    deopt_event_counts: HashMap<u32, usize>,
}

impl DeoptTable {
    pub fn new() -> Self {
        Self {
            entries: HashMap::new(),
            deopt_event_counts: HashMap::new(),
        }
    }

    /// Register metadata for a safepoint.
    pub fn register(&mut self, meta: DeoptMetadata) {
        self.entries.insert(meta.deopt_id, meta);
    }

    /// Retrieve metadata by deopt ID.
    pub fn get(&self, deopt_id: u32) -> Option<&DeoptMetadata> {
        self.entries.get(&deopt_id)
    }

    /// Record a dynamic deoptimization event.
    pub fn record_event(&mut self, deopt_id: u32) {
        *self.deopt_event_counts.entry(deopt_id).or_insert(0) += 1;
    }

    /// Get the total number of times deoptimization was triggered.
    pub fn total_deopts(&self) -> usize {
        self.deopt_event_counts.values().sum()
    }

    /// Clear all registered safepoints and metrics.
    pub fn clear(&mut self) {
        self.entries.clear();
        self.deopt_event_counts.clear();
    }
}

/// Global thread-safe deoptimization registry.
static GLOBAL_DEOPT_TABLE: RwLock<Option<DeoptTable>> = RwLock::new(None);
static GLOBAL_DEOPT_COUNTER: AtomicU64 = AtomicU64::new(0);

/// Initialize the global deopt table.
pub fn init_global_deopt_table() {
    let mut table = GLOBAL_DEOPT_TABLE.write().unwrap();
    *table = Some(DeoptTable::new());
}

/// Register a safepoint in the global deopt table.
pub fn register_global_deopt(meta: DeoptMetadata) {
    let mut guard = GLOBAL_DEOPT_TABLE.write().unwrap();
    if guard.is_none() {
        *guard = Some(DeoptTable::new());
    }
    if let Some(ref mut table) = *guard {
        table.register(meta);
    }
}

/// Reconstruct an unspecialized interpreter call frame mapping variable names to raw values.
pub fn reconstruct_interpreter_frame(
    meta: &DeoptMetadata,
    raw_values: &[i64],
) -> HashMap<String, i64> {
    let mut frame = HashMap::new();
    for (i, (var_name, _, _)) in meta.live_vars.iter().enumerate() {
        let val = if i < raw_values.len() {
            raw_values[i]
        } else {
            0
        };
        frame.insert(var_name.clone(), val);
    }
    frame
}

/// Runtime deoptimization handler called when a speculative type guard fails.
#[no_mangle]
pub extern "C" fn __nl_deopt(deopt_id: u64, _frame_ptr: *const u8) -> i64 {
    GLOBAL_DEOPT_COUNTER.fetch_add(1, Ordering::SeqCst);
    let mut guard = GLOBAL_DEOPT_TABLE.write().unwrap();
    if let Some(ref mut table) = *guard {
        table.record_event(deopt_id as u32);
    }
    0
}

/// Retrieve the total count of runtime deoptimization events.
pub fn get_runtime_deopt_count() -> u64 {
    GLOBAL_DEOPT_COUNTER.load(Ordering::SeqCst)
}

/// Reset runtime deoptimization telemetry.
pub fn reset_runtime_deopt_counter() {
    GLOBAL_DEOPT_COUNTER.store(0, Ordering::SeqCst);
}

/// On-Stack Replacement (OSR) transition slot located in function preambles.
/// Allows dynamic atomic upgrading from unspecialized Tier 0 to specialized Tier 1 native code.
pub struct OsrTransitionSlot {
    /// Atomic function pointer to Tier 1 specialized code (null if unspecialized).
    target_fn_ptr: AtomicPtr<u8>,
    /// Invocation counter tracking execution frequency.
    invocation_count: AtomicU64,
    /// Hotness threshold required to trigger OSR specialization.
    threshold: u64,
}

impl OsrTransitionSlot {
    pub fn new(threshold: u64) -> Self {
        Self {
            target_fn_ptr: AtomicPtr::new(std::ptr::null_mut()),
            invocation_count: AtomicU64::new(0),
            threshold,
        }
    }

    /// Record an entry invocation. Returns true if hot threshold has just been crossed.
    pub fn record_invocation(&self) -> bool {
        let count = self.invocation_count.fetch_add(1, Ordering::Relaxed);
        count + 1 >= self.threshold && self.target_fn_ptr.load(Ordering::Acquire).is_null()
    }

    /// Upgrade the slot with a compiled Tier 1 function pointer.
    pub fn upgrade(&self, specialized_fn: *mut u8) {
        self.target_fn_ptr.store(specialized_fn, Ordering::Release);
    }

    /// Check if specialized code is available for OSR jump.
    pub fn should_osr(&self) -> bool {
        !self.target_fn_ptr.load(Ordering::Acquire).is_null()

\end{lstlisting}
